\begin{abstract}
This technical report presents an end-to-end implementation for emotion analysis in text. The task is to assign scores for eight basic emotions---anger, anticipation, disgust, fear, joy, sadness, surprise, and trust---and to predict the dominant emotion for each input text. The implementation compares three approaches: a classic weighted lexicon baseline based on the NRC Hashtag Emotion Lexicon, a weakly supervised Multinomial Naive Bayes model trained from lexical emotion associations, and a hybrid score-level combination of the two. The system includes command-line prediction for individual texts and CSV files, reproducible experiments, generated metrics, confusion matrices, and visual artifacts. On a small manually constructed evaluation set of 24 texts from social, news, and blog domains, the lexicon baseline obtains the strongest result, with 0.9167 accuracy and 0.9187 macro-F1. The hybrid method improves over Naive Bayes but does not surpass the lexicon baseline on this dataset. The project demonstrates a complete local implementation that satisfies the requirement of combining classic and learned methods while remaining transparent, reproducible, and easy to inspect.
\end{abstract}
